\documentclass{ximera}
%geschreven 13/11/26
\title{Herhaling: arbeid en energie}
\author{Daan Lipkens}

\begin{document}
\begin{abstract}
    Arbeid, kinetische en potentiële zwaarte-energie en de wet van behoud van energie.
\end{abstract}
\maketitle
\label{hoofdstuk:energie}

\section*{Inleiding}
Een kar die je over de vloer duwt, krijgt snelheid. Een steen die je optilt, kan later naar beneden vallen. In beide gevallen wordt energie omgezet: een kracht verricht \emph{arbeid} en verandert zo de energie van een voorwerp.

In het hoofdstuk over potentiële elektrische energie en spanning steunen we op deze begrippen. Daarom zetten we ze in deze appendix kort op een rij:
\begin{itemize}
    \item de arbeid van een constante kracht (sectie~\verwijs[Arbeid]{sec:arbeid}{appendix_energie_arbeid});
    \item kinetische energie en potentiële zwaarte-energie (sectie~\verwijs[Energievormen]{sec:energievormen}{appendix_energie_vormen});
    \item de wet van behoud van energie en de arbeid-energiestelling (sectie~\verwijs[Behoud van energie]{sec:behoud_energie}{appendix_energie_behoud}).
\end{itemize}

\end{document}
